% ======================================================================
\section{The coefficients of \texorpdfstring{$\Omega$}{Omega}: complete expressions}\label{app:coefficients}
% ======================================================================

This appendix collects the coefficients that are not displayed in full in the main text, in momentum space, and all coefficients in transverse coordinate space.

\subsection{Momentum space}

The coefficients $\mf A^{(1)}_2$, $\mf N_2$, $[\mf A^{(3)}_2]_\rho$, $[\mf A^{(3)}_2]_{\rho\rho\rho}$, $\mf A^{(1)}_1$, $\mf B_3$ and $\mf C_1$ are given in Eqs.~\eqref{eq:A1mom}, \eqref{eq:N2mom}, \eqref{eq:A3onerho}, \eqref{eq:A3threerho}, \eqref{eq:A11mom}, \eqref{eq:B3mom} and~\eqref{eq:C1mom}. The two-gluon coefficient is~\cite{Lublinsky:2016meo}
\begin{align}\label{eq:B2mom}
\mf B^{ba}_{2ji}(k,p)=\ &\frac{ig^2f^{abc}\,p^jk^i\rho^c(-\v p-\v k)\sqrt{p^+k^+}}{2\,(p^+\v k^2+k^+\v p^2)}\Big[\frac{1}{k^+\v p^2}-\frac{1}{p^+\v k^2}\Big]\,\Theta(\vee-k^+-p^+)
\notag\\&
+\frac{g^2p^jk^i\,\{\rho^b(-\v p),\rho^a(-\v k)\}}{2\,\v p^2\v k^2\sqrt{p^+k^+}}\,\Theta(\vee-k^+)\,\Theta(\vee-p^+)
\notag\\&
-\frac{ig^2f^{abc}\rho^c(-\v p-\v k)\,(p+k)^l\,(p^mk^+-p^+k^m)}{(\v p+\v k)^2\sqrt{p^+k^+}\,(p^++k^+)\big(\frac{\v p^2}{p^+}+\frac{\v k^2}{k^+}\big)}\notag\\&\quad\times\Big[\frac{\delta_{lm}\delta_{ji}}{p^++k^+}-\frac{\delta_{im}\delta_{jl}}{k^+}-\frac{\delta_{jm}\delta_{il}}{p^+}\Big]\Theta(\vee-k^+-p^+)
\notag\\&
-\frac{ig^2f^{abc}\rho^c(-\v p-\v k)\,(p^+-k^+)\,\delta_{ji}}{2\sqrt{p^+k^+}\,(p^++k^+)^2\big(\frac{\v k^2}{k^+}+\frac{\v p^2}{p^+}\big)}\,\Theta(\vee-k^+-p^+)\,,
\end{align}
for $k^+,p^+>\Lambda$. The first term is the commutator part of the two emissions, the second the anticommutator part, the third the splitting and the fourth the instantaneous emission. The supports in the longitudinal momenta follow from the kinematic restrictions of Section~\ref{sec:lcpt}: in the splitting the parent gluon of momentum $k^++p^+$ is a soft mode, and the same support is used for the instantaneous term; two independently emitted gluons are each soft modes, with no condition on $k^++p^+$.

The part of $\mf A^{(3)}_2$ with two charges is
\begin{align}\label{eq:A3tworho}
\Big[\mf A^{(3)a}_{2i}(k)\Big]_{\rho\rho}=\ &\int d^2p\,\frac{ig^3f^{abc}\{\rho^b(-\v p),\rho^c(-\v k+\v p)\}}{8\sqrt2\,\pi^3\,\v p^2\,\v k^2\sqrt{k^+}}
\Bigg[\log\frac{\Lambda}{k^+}\,\frac{\v k^2p^i-(\v k\cdot\v p)\,k^i}{(\v k-\v p)^2}
-2k^i\log\frac{\vee}{\Lambda}
\notag\\&
-2k^i\log\frac{k^+\v p^2}{\vee\,\v k^2+k^+\v p^2}
-\frac{\v p^2k^i+\v k^2p^i}{(\v k-\v p)^2}\log\frac{\vee-k^+}{\Lambda}-k^i\log\frac{\vee}{k^+}
\notag\\&
+\frac{(2\v p^2-2\v k\cdot\v p+\v k^2)\,k^i+\v k^2p^i}{(\v k-\v p)^2}\log\Big[1+\Big(\frac{\vee}{k^+}-1\Big)\frac{\v k^2}{(\v k-\v p)^2}\Big]\Bigg]\,.
\end{align}
It collects the anticommutator parts of the diagrams with two and three valence vertices; the limit $\Lambda\to0$ has been taken wherever possible.

\subsection{Transverse coordinate space}

In the cross section the coefficients are used in transverse coordinate space, where the eikonal scattering on the target is diagonal. The transverse arguments are defined by the Fourier conventions~\eqref{eq:FTconventions}: $\mf A_2(p^+,-\v x)$ and $\mf A_1(p^+,\v x)$ multiply $a^\dagger(p^+,\v x)$ and $a(p^+,\v x)$, and in general an argument conjugate to the momentum of a creation operator appears with a minus sign. Eq.~\eqref{eq:OmegaFinal} becomes
\begin{align}\label{eq:OmegaX}
\Omega = \ &\mf N
+\int\frac{dp^+}{2\pi}\int_x\Big[\mf A^{(1)b}_{1j}(p^+,\v x)\,a^b_j(p^+,\v x)-\mf A^{(1)\dagger b}_{1j}(p^+,\v x)\,a^{\dagger b}_j(p^+,\v x)
\notag\\&\hspace{2.6cm}
+\mf A^{(3)b}_{1j}(p^+,\v x)\,a^b_j(p^+,\v x)+\mf A^{(3)b}_{2j}(p^+,-\v x)\,a^{\dagger b}_j(p^+,\v x)\Big]
\notag\\&
+\int\frac{dp^+}{2\pi}\frac{dq^+}{2\pi}\int_{x,y}\Big[\mf B^{cb}_{1kj}(p^+,\v x;q^+,\v y)\,a^b_j(p^+,\v x)a^c_k(q^+,\v y)
\notag\\&\hspace{3.2cm}
+\mf B^{cb}_{2kj}(p^+,-\v x;q^+,-\v y)\,a^{\dagger b}_j(p^+,\v x)a^{\dagger c}_k(q^+,\v y)
\notag\\&\hspace{3.2cm}
+\mf B^{cb}_{3kj}(p^+,\v x;q^+,-\v y)\,a^{\dagger c}_k(q^+,\v y)a^b_j(p^+,\v x)\Big]
\notag\\&
+\int\frac{dp^+}{2\pi}\frac{dq^+}{2\pi}\int_{x,y,z}\Big[\mf C^{dcb}_{1lkj}(p^+,\v z-\v y;q^+,\v y-\v x)
\notag\\&\hspace{3.2cm}\times a^{\dagger d}_l(q^+,\v x)a^{\dagger c}_k(p^+-q^+,\v y)a^b_j(p^+,\v z)
\notag\\&\hspace{3.2cm}
-\mf C^{\dagger dcb}_{1lkj}(p^+,\v z-\v y;q^+,\v y-\v x)
\notag\\&\hspace{3.2cm}\times a^{\dagger b}_j(p^+,\v z)a^c_k(p^+-q^+,\v y)a^d_l(q^+,\v x)\Big]+\dots
\end{align}
The coefficients fixed by the diagrammatic calculation are the following. The normalisation and the amplitude of order $g$ are
\begin{align}
\mf N_2&=-\frac{g^2}{8\pi^3}\,\log\frac{\vee}{\Lambda}\int_{x,y,z}K(\v x,\v y,\v z)\,\rho^a(\v x)\,\rho^a(\v y)\,,\label{eq:N2X}\\
\mf A^{(1)a}_{2i}(k^+,-\v z)&=-\mf A^{(1)\dagger a}_{1i}(k^+,\v z)=\frac{ig}{\sqrt2\,\pi\sqrt{k^+}}\int_x\Kv^i(\v x-\v z)\,\rho^a(\v x)\,.\label{eq:A1X}
\end{align}
The amplitude of order $g^3$ for the emission of one gluon has the three parts
\begin{align}\label{eq:A3onerhoX}
\Big[\mf A^{(3)a}_{2i}(k^+,-\v z)\Big]_{\rho}=\ &-\int_x\frac{ig^3N_c\,\rho^a(\v x)\,\Kv^i(\v x-\v z)}{16\sqrt2\,\pi^3\sqrt{k^+}}\Bigg\{\Big[-\frac2\epsilon-2\gamma_E-\log\frac{(\v x-\v z)^2\mu^2_{\MS}}{4}\Big]
\notag\\&\hspace{0.5cm}\times\Big[\frac{11}{3}-2\log\frac{k^+}{\Lambda}+\log\Big(1-\frac{k^+}{\vee}\Big)\Big]
\notag\\&\hspace{0.5cm}
-\frac{67}{9}+\frac{\pi^2}{3}+\log^2\Big(1-\frac{k^+}{\vee}\Big)+\log\Big(1-\frac{k^+}{\vee}\Big)\log\frac{\vee}{k^+}\Bigg\},
\end{align}
where $\gamma_E$ is the Euler constant,
\begin{align}\label{eq:A3tworhoX}
\Big[\mf A^{(3)a}_{2i}(k^+,-\v z)\Big]_{\rho\rho}=\ &\int d^2p\,d^2k\int_{x,y}e^{i\v k\cdot(\v z-\v y)}e^{i\v p\cdot(\v y-\v x)}\,\notag\\&\times\frac{ig^3f^{abc}\{\rho^b(\v x),\rho^c(\v y)\}}{32\sqrt{2k^+}\,\pi^5\,\v p^2\,\v k^2}\,\Big[\ \cdots\ \Big]\,,
\end{align}
with the same square bracket as in Eq.~\eqref{eq:A3tworho} ($\v k$ and $\v p$ are integration variables here), and
\begin{align}\label{eq:A3threerhoX}
\Big[\mf A^{(3)a}_{2i}(k^+,-\v z)\Big]_{\rho\rho\rho}=\ &-\log\frac{\vee}{\Lambda}\int_{w,x,y,u}\rho^a(\v w)\,\rho^b(\v x)\,\rho^b(\v y)\,K(\v x,\v y,\v u)\,\Kv^i(\v w-\v z)\notag\\&\times\frac{ig^3}{8\sqrt2\,\pi^4\sqrt{k^+}}=\mf A^{(1)a}_{2i}(k^+,-\v z)\,\mf N_2\,.
\end{align}
The amplitude for two gluons with longitudinal momenta $k^+$, $p^+$ at the transverse positions $\v w$, $\v z$ is
\begin{align}\label{eq:B2X}
\mf B^{bc}_{2ji}(k^+,-\v w;p^+,-\v z)=\ &\int_x\frac{ig^2f^{cba}\,\rho^a(\v x)\,\sqrt{p^+k^+}}{4\pi^2\big[p^+(\v x-\v z)^2+k^+(\v x-\v w)^2\big](p^++k^+)}
\notag\\&\times\Bigg[\frac{\delta^{ij}\big[(\v x-\v z)^2-(\v x-\v w)^2\big]}{2(\v z-\v w)^2}
\notag\\&\quad
+\frac{p^++k^+}{k^+}\Big(\frac{(x-z)^j(z-w)^i}{(\v z-\v w)^2}-\frac{(x-z)^j(x-w)^i}{2(\v x-\v w)^2}\Big)
\notag\\&\quad
+\frac{p^++k^+}{p^+}\Big(\frac{(x-w)^i(z-w)^j}{(\v z-\v w)^2}+\frac{(x-z)^j(x-w)^i}{2(\v x-\v z)^2}\Big)\Bigg]
\notag\\&\quad\times\Theta(\vee-k^+-p^+)
\notag\\&
-\int_{x,y}\frac{g^2\,\{\rho^c(\v y),\rho^b(\v x)\}\,\Kv^i(\v y-\v w)\,\Kv^j(\v x-\v z)}{8\pi^2\sqrt{p^+k^+}}
\notag\\&\quad\times\Theta(\vee-k^+)\,\Theta(\vee-p^+)\,.
\end{align}
The amplitude for the transition of a gluon $(k^+,\v z,a,i)$ into a gluon $(p^+,\v u,b,j)$ on the valence charges is
\begin{align}\label{eq:B3X}
\mf B^{ba}_{3ji}(k^+,\v z;p^+,-\v u)=\ &\frac{ig^2f^{abc}}{(2\pi)^2}\int_x\frac{\sqrt{p^+k^+}\,(p^++k^+)}{(k^+-p^+)^2}\,
\notag\\&\qquad\times\frac{\rho^c(\v x)\,\delta_{ij}}{p^+(\v x-\v u)^2-k^+(\v x-\v z)^2}\,\big[\Theta_>+\Theta_<\big]
\notag\\&
+\frac{ig^2f^{abc}}{2\pi^2k^+}\int_x\frac{\sqrt{k^+p^+}}{k^+-p^+}\,\rho^c(\v x)
\notag\\&\times\frac{\big(k^+(x-z)^m-p^+(x-u)^m\big)\big((x-z)^l-(x-u)^l\big)}{k^+(\v x-\v z)^2-p^+(\v x-\v u)^2}
\notag\\&\qquad\times\Big(\delta_{il}\delta_{jm}-\frac{k^+}{p^+}\delta_{jl}\delta_{im}-\frac{k^+}{k^+-p^+}\delta_{ij}\delta_{lm}\Big)
\notag\\&\qquad\times\Bigg[-\frac{\Theta_>}{(\v u-\v z)^2}-\frac{p^+k^+\,\Theta_<}{\big(k^+(\v x-\v z)-p^+(\v x-\v u)\big)^2}\Bigg]
\notag\\&
-\frac{g^2}{4\pi^2\sqrt{p^+k^+}}\int_{x,y}\Kv^i(\v x-\v z)\,\Kv^j(\v y-\v u)\,\{\rho^a(\v x),\rho^b(\v y)\}
\notag\\&
-\frac{ig^2f^{abc}}{4\pi^2\sqrt{p^+k^+}}\int_x\Kv^i(\v x-\v z)\,\Kv^j(\v x-\v u)
\notag\\&\qquad\times\frac{k^+(\v x-\v z)^2+p^+(\v x-\v u)^2}{p^+(\v x-\v u)^2-k^+(\v x-\v z)^2}\,\rho^c(\v x)\,,
\end{align}
with principal values for the denominators $p^+(\v x-\v u)^2-k^+(\v x-\v z)^2$. It follows from Eqs.~\eqref{eq:psi_instX}, \eqref{eq:psi1X}, \eqref{eq:psi2X}, \eqref{eq:psi_antiX} and~\eqref{eq:psi_commX} with $\v X'=\v x-\v z$, $\v X=\v x-\v u$, $\xi=p^+/k^+$ and the factor $2\pi/k^+$ between the coefficient of $\ket{g^b_j(\xi k^+,\v u)}$ under $\int d\xi\int_u$ and $\mf B_3$. The amplitude for the splitting of a gluon $(k^+,a,i)$ into the gluons $(p^+,b,j)$ and $(k^+-p^+,c,k)$ is
\begin{align}\label{eq:C1X}
\mf C^{bca}_{1jki}(k^+,\v z;p^+,\v x)=\ &-\frac{gf^{abc}}{2\pi\sqrt{2k^+}}\,\frac{\sqrt{p^+(k^+-p^+)}}{k^+}\,\delta^{(2)}\Big(\v z+\frac{p^+}{k^+}\v x\Big)\,\Kv^m(\v x)
\notag\\&\times\Big[\delta_{im}\delta_{jk}-\frac{k^+}{p^+}\delta_{jm}\delta_{ik}-\frac{k^+}{k^+-p^+}\delta_{km}\delta_{ij}\Big]\Theta(k^+-p^+-\Lambda)\,,
\end{align}
where the arguments are those of Eq.~\eqref{eq:OmegaX}: $\v z$ is the position of the incoming gluon relative to the outgoing gluon $(c,k)$, and $\v x$ the position of the latter relative to the outgoing gluon $(b,j)$. The $\delta$-function expresses that the incoming gluon sits at the centre of the two outgoing gluons weighted with their longitudinal momenta.

The coefficients fixed by unitarity are, in coordinate space,
\begin{align}\label{eq:unitX}
\mf C^{\dagger dcb}_{2lkj}(p^+,\v z;q^+,\v x)&=-\mf C^{dcb}_{1lkj}(p^+,-\v z;q^+,-\v x)\,,\notag\\
\mf B^{\dagger cb}_{1kj}(p^+,\v x;q^+,\v y)&=-\mf B^{cb}_{2kj}(p^+,-\v x;q^+,-\v y)\notag\\
&\quad+\frac12\Big[\mf A^{(1)b}_{2j}(p^+,-\v x)\,\mf A^{(1)c}_{2k}(q^+,-\v y)\notag\\&\qquad\quad+\mf A^{(1)c}_{2k}(q^+,-\v y)\,\mf A^{(1)b}_{2j}(p^+,-\v x)\Big]
\notag\\&\quad
+\int_z\mf C^{cbd}_{1kjl}(q^++p^+,\v z-\v x;q^+,\v x-\v y)\,\mf A^{(1)d}_{2l}(q^++p^+,-\v z)\,,\notag\\
\mf A^{(3)\dagger b}_{1j}(p^+,\v x)&=-\mf A^{(3)b}_{2j}(p^+,-\v x)-\big[\mf N_2,\mf A^{(1)b}_{2j}(p^+,-\v x)\big]
\notag\\&\quad+\int\frac{dq^+}{2\pi}\int_y\Big[2\,\mf A^{(1)c}_{1k}(q^+,\v y)\,\mf B^{cb}_{2kj}(p^+,-\v x;q^+,-\v y)
\notag\\&\qquad
-\mf B^{\dagger cb}_{3kj}(p^+,\v x;q^+,-\v y)\,\mf A^{(1)c}_{2k}(q^+,-\v y)
\notag\\&\qquad-2\int_z\mf C^{\dagger dcb}_{1lkj}(p^+,\v x-\v z;q^+,\v z-\v y)\,\mf B^{dc}_{2lk}(p^+-q^+,-\v z;q^+,-\v y)\Big]\,.
\end{align}

\NOTE{Source of this appendix: the coordinate-space expressions~\eqref{eq:A3onerhoX}--\eqref{eq:C1X} and~\eqref{eq:unitX} are taken from the checked calculation file; the transforms from Eq.~\eqref{eq:B2mom} to Eq.~\eqref{eq:B2X} and from Eq.~\eqref{eq:C1mom} to Eq.~\eqref{eq:C1X}, and the transforms~\eqref{eq:FT_i4} and~\eqref{eq:FT_i5}, were not re-derived here.}

